\subsection{High-dimensional experiments with QuadraSHAP}
\label{sec:high-dimensional-quadrashap}

To assess whether Shapley attribution remains practical when the number of
features greatly exceeds the number of patients, we consider two genomic
survival datasets: GSE24080 multiple-myeloma gene expression
\cite{geo_gse24080} and TCGA lower-grade glioma (LGG) DNA methylation
\cite{tcga_lgg_methylation}.
Genome-wide assays measure many molecular variables per patient, whereas
the cohorts with matched survival follow-up contain only hundreds of patients.
After eligibility filtering, GSE24080 contains 553 patients and 54,675
expression probe sets, and TCGA LGG contains 511 patients and 396,065
methylation sites. These are probe- or site-level features, retained without
gene aggregation. We use the supplied 339/214 training/validation split
for GSE24080 and a 383/128 split stratified by event status for TCGA LGG
(seed 42), giving an extreme small-$n$, large-$d$ setting.

We model overall survival using a Cox proportional-hazards model, which
incorporates right-censored observations: patients whose death has not been
observed by the end of follow-up. Its relative hazard factors across features,
making QuadraSHAP applicable directly:
\[
h(t\mid x)=h_0(t)f(x),
\qquad
f(x)=\exp(\beta^\top x)=\prod_{j=1}^{d}\exp(\beta_jx_j).
\]
To stabilize estimation when $d\gg n$, we shrink the Cox coefficients
towards zero using a ridge penalty, with the implementation's penalty
parameter fixed a priori at $\alpha=0.01d$. We use the Breslow approximation
for patients with equal recorded event times. Missing values are imputed using
training medians, and centering and scaling use training statistics.
All retained features have nonzero fitted coefficients.

We explain the relative hazard using 30 training-background patients,
sampled without replacement with seed 42 and held fixed across methods.
Specifically, absent features are replaced by the corresponding entries
of each background row and predictions are averaged:
\[
v_x(S)=\frac{1}{30}\sum_{b=1}^{30}
f\bigl(x_S,z^{(b)}_{\overline S}\bigr).
\]
The resulting game is a weighted sum of product games. We explain the
first five held-out patients in each saved split with the degree-exact rule
$m_q=\lceil d/2\rceil$ and with patient-specific node counts selected
to certify absolute quadrature tolerances
$\varepsilon\in\{10^{-1},10^{-3},10^{-6}\}$.
All methods use JAX-Metal FP32 on an Apple M4 Pro, an 8\,GiB memory-planning
budget, and one background row per block. We record one synchronized
explanation call per patient and method after warming up each distinct
dataset/node-count configuration. Reported times include automatic node
selection, data transfers, and host accumulation, but exclude model fitting,
warm-up, and cached quadrature-rule construction.

Table~\ref{tab:cox-high-dimensional} shows that degree-exact computation
remains manageable even at these dimensions: 27,338 and 198,033 nodes
require an average of 1.17\,s and 66.38\,s per patient, respectively.
For latency-sensitive use, the certified node budget reduces computation
to 85--89\,ms for GSE24080 and 374--385\,ms for TCGA LGG.
At $\varepsilon=10^{-1}$, this corresponds to approximately
$13.8\times$ and $177.2\times$ speedups, with all five patients in each
cohort meeting the requested discrepancy threshold against the GPU
degree-exact result. These subsecond latencies support interactive,
on-demand Shapley explanations once the model and quadrature rules are ready.
Caching matters: constructing the large exact rules previously required
10.97\,s and 561.92\,s on CPU, respectively, whereas approximate-rule
construction took less than 1\,ms.

The numerical comparison also identifies the scope of the certificate.
For patient $p$, let $\phi_p$ denote the exact Shapley vector of the fixed
background game and $\phi_{m_p,p}$ its quadrature approximation in real
arithmetic. The selected node count satisfies
\[
\|\phi_{m_p,p}-\phi_p\|_\infty\le C_p\le\varepsilon,
\qquad C_{\max}=\max_{p=1,\ldots,5}C_p.
\]
The measured quantity instead compares floating-point outputs,
$\Delta_{\mathrm{GPU}}=
\max_p\|\widehat\phi^{\mathrm{GPU}}_{m_p,p}
-\widehat\phi^{\mathrm{GPU}}_{\mathrm{exact},p}\|_\infty$.
Although every selected rule satisfies its quadrature bound,
the measured $10^{-3}$ threshold is met for only 4/5 GSE24080 patients
and 2/5 TCGA LGG patients, and the $10^{-6}$ threshold is met for none.
Moreover, GPU degree-exact outputs themselves differ from an independent,
tightly converged FP64 reference by up to $1.87\times10^{-3}$ and
$3.32\times10^{-2}$, respectively, so agreement with that GPU reference
does not establish exact numerical accuracy.
These experiments therefore demonstrate rapid, quadrature-certified
approximation, while stringent end-to-end GPU error guarantees require
additional control of floating-point error. The certificate applies to
the chosen 30-row background game and does not bound uncertainty from
approximating the background distribution.
